% ===========================================================================
% chapters/01_introduction.tex
% ===========================================================================
\chapter{Introduction}
\label{ch:introduction}

We can now predict, before training begins, how much a large language model's loss will fall
as its parameters, data, and compute are scaled, through the neural scaling laws that have
become a central tool for planning large-scale training~\cite{kaplan2020scaling,hoffmann2022chinchilla}.
A scaling law, however, describes \emph{how much} a model will improve, not \emph{what} it is
learning or \emph{when} during training each capability and internal structure takes shape.
That internal trajectory---the order and manner in which a model's representations are laid
down beneath a smoothly falling loss---is the subject of training-dynamics research, and it is
the subject of this thesis. Understanding it matters: it would let us reason about when a
capability appears, design data and curricula that bring useful structure forward, anticipate
emergent behaviour, and---as this thesis emphasizes---reason about \emph{when} during training
the model is most amenable to being shaped.

Viewed at the level of internal structure rather than aggregate loss, training is far less
smooth than the loss curve suggests. Generalization can appear abruptly long after the loss
has plateaued~\cite{power2022grokking}; such plateaus often conceal a continuously improving
internal quantity, so that the eventual breakthrough is the surfacing of hidden progress
rather than a sudden change~\cite{barak2022hidden,nanda2023progress}; and capabilities have
been reported to appear suddenly with scale, although part of this apparent abruptness is an
artifact of harsh evaluation metrics rather than of the underlying
competence~\cite{wei2022emergent,schaeffer2023mirage}. Analyses of the internal geometry of
real pre-training runs reinforce the picture: representations pass through a sequence of
qualitatively distinct phases over training~\cite{representation_geometry_phases}. These
findings establish that pre-training has a rich \emph{developmental structure}---phases,
transitions, and hidden accumulation---beneath the monotone loss.

Most of this work, however, describes the trajectory of a \emph{single} capability or of
\emph{global} geometry, and treats the phases as something to characterize rather than as
something with consequences. A different and largely open question is whether the phases are
\emph{functionally} different: whether the model is differently malleable at different points
in training, such that the same input leaves a different lasting mark depending on when it
arrives. This is the defining idea of a \emph{critical} or \emph{sensitive period}---a window
during which a developing system is maximally and durably shaped by its input---a concept
central to biological learning and demonstrated in artificial networks, where perturbations
early in training leave lasting effects that the identical perturbation applied later does
not~\cite{achille2019critical}. Whether large language model pre-training has analogous
windows, and whether they can be located from the model's own internal indicators, is the gap
this thesis addresses.

The question has an immediate and consequential application. Contemporary alignment follows a
\emph{post-hoc} paradigm: models pre-train on massive, largely unfiltered corpora, after which
safety is imposed by fine-tuning. There is mounting evidence that this produces \emph{fragile}
alignment---small quantities of poisoned data can implant durable backdoors~\cite{poisoning2025},
and narrow fine-tuning can induce broad ``emergent misalignment''~\cite{betley2025emergent}.
If pre-training contains a developmental window during which signal is more deeply and durably
internalized, then \emph{when} alignment-relevant signal is introduced---not only how much of
it the model sees---would matter, with direct implications for training recipes. We are careful
throughout to treat this as motivation and as a downstream implication, gated on the
developmental result, rather than as a claim the thesis assumes.

This project originates in exactly this question. We began by tracking weight-spectral
indicators---effective rank, stable rank, spectral and Frobenius norms, the count of
eigenvalues separating from the random-initialization bulk, eigenspectrum decay, and the
stability of singular vectors across checkpoints---across the public Pythia checkpoint
suite~\cite{biderman2023pythia}, layer by layer and module by module, asking whether
pre-training has separable phases. The first completed experiment in this thesis identifies a
reproducible early weight-space reorganisation across Pythia-70M, 160M, 410M, and 1B,
concentrated in the first few thousand optimisation steps and strongest between approximately
512 and 2000 steps. This observational result defines the candidate window; the rest of the
thesis asks whether that window is functionally meaningful and whether it marks a period of
heightened durable plasticity. The research questions are therefore:

\begin{researchquestion}[Phases]
\label{rq:phases}
Does language model pre-training pass through identifiable developmental phases, detectable
from internal weight-spectral indicators, and are those phases \emph{functionally meaningful}
in the sense of coinciding with changes in what the model can do?
\end{researchquestion}

\begin{researchquestion}[Critical period]
\label{rq:critical}
Is there a window during pre-training in which a signal injected into the model is more durably
internalized---surviving subsequent training and resisting later degradation---than the same
signal injected later, and does this window's boundary coincide with the indicator-defined
consolidation transition rather than merely with the learning-rate schedule?
\end{researchquestion}

\begin{researchquestion}[Structure and generality]
\label{rq:structure}
Is any such window global or module-specific; does it replicate across distinct signal types
(making it a property of the developmental stage rather than of the signal); and, as the
structured case, is alignment-style signal injected within the window more robust to
degradation and poisoning-style attack than post-hoc alignment?
\end{researchquestion}

The thesis pursues these questions in a deliberately tiered way. The developmental phases are
located from internal indicators measured across real Pythia checkpoints, independently of any
injected signal; the critical-period claim is then tested \emph{causally} by using the
checkpoints as injection points---introducing a defined signal at different training stages and
measuring how durably it is retained. Because the window is defined from weight geometry and
durability is measured behaviourally, the central test compares two independent measurements,
which is what protects it from circularity. Controlled toy models, in which a phase transition
is known by construction, are used to calibrate the interpretation of the indicators and to
supply mechanistic vignettes; they are not relied upon to establish the claim at scale.

Contingent on the results, the contributions are as follows. First, a characterization of
pre-training as a phased developmental process, established from convergent weight-spectral
indicators across model sizes, including the global phase structure and its module- and
layer-wise heterogeneity. Second, a test of whether these phases are functionally meaningful by
aligning them with loss and downstream-capability trajectories. Third, the central
contribution: a causal test of the critical-period hypothesis via stage-resolved signal
injection, with the discriminating prediction that durability exhibits a \emph{break} at the
indicator-defined consolidation boundary rather than a trivial monotone decay, evaluated against
learning-rate and cumulative-exposure baselines. Fourth, an analysis of whether the window is
module-specific and whether it generalizes across signal types, with alignment durability as the
applied case connecting the developmental finding to the fragility of post-hoc alignment. The
design ties each outcome to a defensible claim: a clean window that closes at the consolidation
boundary would establish a critical period; a monotone decay would constitute an honest negative
that the effect is the trivial ``early matters'' phenomenon; and a module-resolved window would
refine the account into a statement about \emph{where} as well as \emph{when} signal is durably
acquired.

The remainder of the thesis proceeds as follows. \Cref{ch:background} reviews scaling and
emergence, the hypotheses on the temporal structure of training (quantization, hidden progress,
implicit curriculum), critical learning periods, the geometry of pre-training, the fragility of
post-hoc alignment, and the spectral tools used to locate phases, and positions the contribution.
\Cref{ch:methodology} specifies the indicators and phase-identification procedure, the
intervention and durability measures, the role of toy models, and the justification design that
makes each result defensible. \Cref{ch:results} reports the completed E1 phase-identification
result and then lays out the functional-grounding and intervention experiments that test the
stronger behavioural claims. \Cref{ch:conclusion} synthesizes the claims, their boundaries, and
their implications.
